\heading{经济学原理（II）-样题}
\noteinfo{本文件依据你上传的“期中考试（样例）”与“样例参考答案”整理为可直接放入模板的 TeX 版本。原样题中的若干图片题在这里改写为文字自洽的形式，但知识点、题型结构与答案要点保持一致。}

\instructions{满分 100 分，建议答题时间 120 分钟。对于计算题，请写出关键中间步骤。}

\textbf{一、单项选择题}

\begin{question}
原样题第 1 题配有一条“黄油—大炮”的向外凸生产可能性边界：A 位于纵轴上端附近，B 位于边界上部，C 位于曲线内侧，D 位于靠近横轴的边界点。问：在这些点中，生产大炮的机会成本最低的是哪一点？
\begin{choiceoptions}[2]
  \optionitem{A}
  \optionitem{B}
  \optionitem{C}
  \optionitem{D}
\end{choiceoptions}
\begin{solution}
沿着向外凸的生产可能性边界，从左上向右下移动时，生产更多大炮的边际机会成本递增。C 点位于曲线内侧，意味着资源未充分利用；若比较“再多生产一点大炮所需放弃的黄油”，样题标准答案给出 C。核心考点是识别“机会成本最低”对应曲线较平缓、或在非效率点存在改进空间。\anstext{C}
\end{solution}
\end{question}

\begin{question}
下面哪一辆轿车在被交易时\textbf{不影响} 2025 年中国的 GDP？
\begin{choiceoptions}[1]
  \optionitem{一辆 2025 款全新蔚来 ES8 被一家中国国企采购，用于接送领导}
  \optionitem{一辆仓库里全新的 2024 款蔚来 ES8 被一家杭州盈利性中间商卖给一家美国私企}
  \optionitem{一辆 2025 款全新 Toyota 轿车在日本生产并卖给正在当地旅游的小李，随后小李将车开回上海并缴纳关税}
  \optionitem{一辆 2025 款南京生产的上汽大众 Polo Plus 被一对夫妇低价转手给隔壁老王}
\end{choiceoptions}
\begin{solution}
A、B、C 都对应当期新生产商品或当期服务，可能进入 2025 年 GDP；D 是二手车转手，车辆本身不是 2025 年新生产，不计入当期 GDP。\anstext{D}
\end{solution}
\end{question}

\begin{question}
下面哪位经济学家发表了《通论》，引发了经济学研究范式的革命？
\begin{choiceoptions}[2]
  \optionitem{亚当·斯密}
  \optionitem{约翰·梅纳德·凯恩斯}
  \optionitem{西蒙·库兹涅茨}
  \optionitem{米尔顿·弗里德曼}
\end{choiceoptions}
\begin{solution}
《就业、利息和货币通论》由凯恩斯发表，是宏观经济学史上的关键著作。\anstext{B}
\end{solution}
\end{question}

\begin{question}
下面哪一个说法是不恰当的？
\begin{choiceoptions}[1]
  \optionitem{学习经济学的一个好处是你最终会理解在股市中如何赚钱}
  \optionitem{学习经济学可以帮助你思考为了防止致命病毒扩散的经济封锁政策面临的权衡取舍}
  \optionitem{学习经济学的一个好处是你将培养做出更好决策的能力}
  \optionitem{在课表中某个宝贵时间段选择一门经济学课程，其机会成本是放弃其他同时段课程}
\end{choiceoptions}
\begin{solution}
经济学训练的是分析框架，不保证你一定在股市赚钱。其余说法都体现了经济学关于权衡、选择与机会成本的基本思想。\anstext{A}
\end{solution}
\end{question}

\begin{question}
泰德曾经是一名高中数学老师，他辞职后重返大学攻读计算机编程；毕业后被一家软件公司高薪聘用。这个例子最能体现下面哪个理论？
\begin{choiceoptions}[2]
  \optionitem{补偿性工资差异}
  \optionitem{效率工资}
  \optionitem{信号}
  \optionitem{人力资本}
\end{choiceoptions}
\begin{solution}
重新接受教育提高了生产率并带来更高工资，最直接体现人力资本理论。\anstext{D}
\end{solution}
\end{question}

\begin{question}
我们很难衡量歧视如何影响人们在劳动力市场上的结果，主要是因为
\begin{choiceoptions}[2]
  \optionitem{工资数据是关键但不易获取}
  \optionitem{企业通常会谎报工资以掩饰歧视}
  \optionitem{人们有不同的个人特征，且就职于不同岗位}
  \optionitem{最低工资法无差别地应用于不同人群}
\end{choiceoptions}
\begin{solution}
要识别歧视效应，必须把教育、经验、行业、岗位、地区等差异控制住，而现实中个体异质性很多，因此识别困难。\anstext{C}
\end{solution}
\end{question}

\begin{question}
约翰·罗尔斯关于“无知者之幕”背后的“原位”的思想实验，本意是引起人们注意
\begin{choiceoptions}[2]
  \optionitem{大多数穷人不知道如何找到更好的工作}
  \optionitem{每个人出生后落入的人生境遇大多数由运气决定}
  \optionitem{富人有太多钱了，不知道怎么花}
  \optionitem{只有平等机会才一定有效率}
\end{choiceoptions}
\begin{solution}
罗尔斯希望人们在不知道自己会出生在社会哪个位置时来思考制度设计，从而意识到初始境遇很大程度带有运气成分。\anstext{B}
\end{solution}
\end{question}

\begin{question}
根据摩尔定律，集成电路上晶体管数目约每两年增加一倍。根据增长率近似公式，这意味着晶体管数目的年增长率大约为
\begin{choiceoptions}[2]
  \optionitem{41\%}
  \optionitem{100\%}
  \optionitem{50\%}
  \optionitem{30\%}
\end{choiceoptions}
\begin{solution}
若两年翻一倍，则年增长率约满足 $(1+g)^2=2$，所以 $g\approx \sqrt2-1\approx 41\%$。\anstext{A}
\end{solution}
\end{question}

\begin{question}
经济增长的稳态均衡意味着什么不随时间发生改变？
\begin{choiceoptions}[2]
  \optionitem{不平等}
  \optionitem{贫困率}
  \optionitem{人均 GDP}
  \optionitem{人均实物资本存量}
\end{choiceoptions}
\begin{solution}
索罗模型中的稳态指人均资本不再变化；人均 GDP 在无技术进步的稳态下也不再增长，但样题标准答案对应的是“人均实物资本存量”。\anstext{D}
\end{solution}
\end{question}

\begin{question}
约翰和小芳结婚并居住在中国。约翰是加拿大公民，在家照顾两个孩子，每个孩子照料价值各 20 万元；他夜里还为一个日本网站编写游戏程序，每年赚取 40 万元。小芳是中国公民、律师，因为业务常去美国出差，并赚取 200 万元年收入。请问，二人加起来对中国 GNP 的贡献是多少？
\begin{choiceoptions}[2]
  \optionitem{40 万}
  \optionitem{80 万}
  \optionitem{200 万}
  \optionitem{280 万}
\end{choiceoptions}
\begin{solution}
GNP 按居民口径计算。约翰不是中国居民（按样题给定口径），其家庭照料也不是市场交易；他为日本网站编程收入不计入中国 GNP。小芳是中国居民，其 200 万收入计入中国 GNP，因此总贡献为 200 万。\anstext{C}
\end{solution}
\end{question}

\textbf{二、简答题}

\begin{question}
请简单解释纵向公平与横向公平。这两个对公平的理解是由什么原则引申出来的？学界对这两个理解的争论是什么？
\begin{solution}
纵向公平是指支付能力更强的人应纳更多税；横向公平是指支付能力相似的人应承担相似税负。二者都引申自\textbf{支付能力原则}。争论点在于：纵向公平里，“多纳多少”才算公平并无唯一答案；横向公平里，“何谓相似支付能力”也并不容易界定，例如是否要把家庭责任、健康状况和非货币福利计入。\end{solution}
\end{question}

\begin{question}
当存在买方垄断时，为什么雇主能够以低于边际产品价值的工资招到所需的工人？
\begin{solution}
买方垄断雇主面对向上倾斜的劳动供给曲线。若想多雇一个人，往往必须提高所有工人的工资，因此\textbf{边际雇佣成本}高于工资本身。企业按照“边际雇佣成本等于劳动边际产品价值”决定雇佣量，于是均衡工资会低于边际产品价值。这说明买方垄断会压低工资与就业。\end{solution}
\end{question}

\begin{question}
请列举四个导致我国女性在劳动力市场中获得的工资逐年上涨的可能原因。
\begin{solution}
可答的方向很多，写出四点即可。例如：\textbf{(1)} 女性受教育程度上升，人力资本提高；\textbf{(2)} 对女性的歧视逐步下降；\textbf{(3)} 科技进步与新产业扩大了女性更容易进入的高薪岗位；\textbf{(4)} 女性劳动参与和职业晋升机会增加；\textbf{(5)} 最低工资或工会力量改善了弱势女性工资；\textbf{(6)} 更多女性愿意进入高压力高回报职业，从而获得补偿性工资差异。\end{solution}
\end{question}

\begin{question}
请比较功利主义与自由主义。这两种政治哲学会导致绝对平等吗？为什么？
\begin{solution}
功利主义主张最大化社会总效用；自由主义（此处按罗尔斯意义）强调在“无知之幕”后选择公正制度，更重视最不利者处境。二者都\textbf{不必然导向绝对平等}。功利主义会考虑边际效用递减，因此支持一定再分配，但也会担心激励扭曲导致总财富下降；罗尔斯自由主义更接近平等，因为它更强调最不利者利益，但同样不会忽视激励和效率，故也不必然要求结果完全相同。\end{solution}
\end{question}

\textbf{三、计算和分析题}

\begin{question}
请将下表补充完整，并通过至少三种 GDP 计算方法说明，不同方法算出的 GDP 在理论上应当相等。公司 1、2、3 分别代表产业链上、中、下游企业。

\begin{center}
\begin{tabular}{c|cccc}
 & 公司1 & 公司2 & 公司3 & 经济整体 \\\hline
销售收入 & 1500 & 2000 & \blank[1.2cm] & -- \\
原料投入 & 0 & \blank[1.2cm] & \blank[1.2cm] & -- \\
劳动报酬 & 500 & 100 & \blank[1.2cm] & 9000 \\
地租 & 100 & 75 & 200 & \blank[1.2cm] \\
利息 & \blank[1.2cm] & 90 & 500 & \blank[1.2cm] \\
利润 & 65 & 235 & \blank[1.2cm] & 3200 \\
增加值 & \blank[1.2cm] & \blank[1.2cm] & 12000 & \blank[1.2cm]
\end{tabular}
\end{center}
\begin{solution}
由公司 1 可得增加值 $=1500-0=1500$，故利息 $=1500-500-100-65=835$。公司 2 的原料投入应为来自公司 1 的 1500，因此增加值 $=2000-1500=500$。公司 3 的原料投入应为来自公司 2 的 2000，故销售收入 $=12000+2000=14000$，劳动报酬 $=12000-200-500-2900=8400$，利润 $=3200-65-235=2900$。经济整体地租 $=100+75+200=375$，利息 $=835+90+500=1425$，GDP $=1500+500+12000=14000$。

支出法（最终产品法）：\ $GDP=14000$。

收入法：\
\[
GDP=9000+375+1425+3200=14000.
\]

增加值法：\
\[
GDP=1500+500+12000=14000.
\]
\ansmath{GDP=14000}
\end{solution}
\end{question}

\begin{question}
下表给出 2021 年美国家庭五等份收入份额：最低组 3.6\%，第二组 9.0\%，第三组 14.6\%，第四组 22.6\%，最高组 50.1\%；同时已知收入最高 5\% 的份额是 22.1\%。
\begin{enumerate}[label=(\alph*),leftmargin=2.5em]
  \item 请基于这些分组数据说明洛伦兹曲线应如何绘制；
  \item 计算基尼系数，并说明与真实基尼系数相比，你的估计为什么会有误差。
\end{enumerate}
\begin{solution}
横轴取人口累计占比，纵轴取收入累计占比。用五等份并进一步把最高 20\% 拆成 80--95\% 和 95--100\% 两段后，可得到点：
\[
(0,0),\ (0.2,0.0362),\ (0.4,0.1264),\ (0.6,0.2726),\ (0.8,0.4988),\ (0.95,0.7790),\ (1,1).
\]
其中小数略按样题做了凑整处理。

洛伦兹曲线下方面积近似为
\[
A_L\approx 0.27723.
\]
因此基尼系数
\[
G=1-2A_L=1-2\times 0.27723=0.44554.
\]
\ansmath{G\approx 0.4455}

误差的来源在于：我们把每一组内部家庭都近似看作“收入相同”，真实分布显然更连续、更厚尾；不过把最高收入组进一步拆出最高 5\% 后，已经比只用五等份更准确，因为它更好刻画了高收入尾部。\end{solution}
\end{question}

\begin{question}
一个水果国度中，三种水果在 2023--2025 年的产量与价格如下：
\begin{center}
\begin{tabular}{c|cc|cc|cc}
 & \multicolumn{2}{c|}{2023年} & \multicolumn{2}{c|}{2024年} & \multicolumn{2}{c}{2025年} \\
水果 & 产量 & 价格 & 产量 & 价格 & 产量 & 价格 \\\hline
苹果 & 30 & 1 & 40 & 1 & 45 & 4 \\
橙子 & 90 & 2 & 95 & 3 & 90 & 5 \\
西柚 & 60 & 1 & 65 & 2 & 70 & 3
\end{tabular}
\end{center}
\begin{enumerate}[label=(\alph*),leftmargin=2.5em]
  \item 求 2024 至 2025 年名义 GDP 增长率；
  \item 以 2024 年为基期，求 2023 年真实 GDP；
  \item 求 2024 年的费雪环比经济指数；
  \item 若代表性居民每年消费 3 个苹果、9 个橙子和 6 个西柚，以 2023 年为基期，求 2025 年 CPI。
\end{enumerate}
\begin{solution}
2024 年名义 GDP 为
\[
40\times1+95\times3+65\times2=455.
\]
2025 年名义 GDP 为
\[
45\times4+90\times5+70\times3=840.
\]
故名义 GDP 增长率为
\[
\frac{840-455}{455}\times 100\%=84.6\%.
\]

以 2024 年价格计算 2023 年真实 GDP：
\[
30\times1+90\times3+60\times2=420.
\]

2024 年相对 2023 年的费雪环比经济指数为
\[
F_{2024,2023}=\sqrt{\frac{40\times1+95\times2+65\times1}{30\times1+90\times2+60\times1}\times
\frac{40\times1+95\times3+65\times2}{30\times1+90\times3+60\times2}}\times 100
\approx 108.8.
\]

CPI 以 2023 年为基期：
\[
\text{2023篮子成本}=3\times1+9\times2+6\times1=27,
\]
\[
\text{2025篮子成本}=3\times4+9\times5+6\times3=75,
\]
所以
\[
CPI_{2025}=\frac{75}{27}\times 100\approx 277.8.
\]
\ansmath{\text{名义 GDP 增长率}=84.6\%,\quad RGDP_{2023|2024}=420,}
\ansmath{F_{2024,2023}\approx 108.8,\quad CPI_{2025|2023}\approx 277.8}
\end{solution}
\end{question}

\begin{question}
某国因疫情冲击偏离稳态。已知 2024 年人口为 2 亿，资本存量为 4.5 万亿，人口增长率为 5\%，消费占 GDP 比重为 70\%，总生产函数为
\[
Y=K^{0.5}L^{0.5}.
\]
假设没有折旧和技术进步。
\begin{enumerate}[label=(\alph*),leftmargin=2.5em]
  \item 补全 2024 和 2025 年的人口、资本、GDP、储蓄（投资）、维持资本密度所需投资、资本密度、人均 GDP 和人均消费；
  \item 若平均收敛速度取第 1 年的 36\%，估计 2068 年人均 GDP 距离稳态还差多少；
  \item 政府应鼓励储蓄还是消费？黄金规则储蓄率目标是多少？
\end{enumerate}
\begin{solution}
先算 2024 年：
\[
Y=\sqrt{4.5\times 2}=3,\qquad s=0.3,
\]
故储蓄（投资）
\[
I=sY=0.9.
\]
维持资本密度不变所需投资为
\[
nK=0.05\times 4.5=0.225.
\]
资本密度为
\[
k=\frac{4.5}{2}=2.25,
\]
人均 GDP 为
\[
y=\frac{3}{2}=1.5,
\]
人均消费为
\[
(1-s)y=0.7\times 1.5=1.05.
\]

到 2025 年，人口变为 $2.1$ 亿，资本变为
\[
K'=4.5+0.9=5.4.
\]
故
\[
Y'=\sqrt{5.4\times 2.1}=3.3675,
\]
\[
I'=0.3\times 3.3675=1.0102,
\]
\[
\text{维持资本密度投资}=0.05\times 5.4=0.27,
\]
\[
k'=\frac{5.4}{2.1}=2.5714,
\]
\[
y'=\frac{3.3675}{2.1}=1.6036,
\]
\[
\text{人均消费}=0.7\times 1.6036=1.1225.
\]

稳态满足
\[
0.3k_e^{0.5}=0.05k_e\Rightarrow k_e=36,
\]
故稳态人均产出
\[
y_e=\sqrt{36}=6.
\]
第 1 年人均 GDP 增长率近似
\[
\frac{1.6036-1.5}{1.5}\approx 0.069.
\]
按题意用其中 36\% 作为平均收敛速度近似，44 年后距稳态还差
\[
6-1.5\times (1+0.069\times 0.36)^{44}\approx 1.6\text{ 万/人}.
\]

黄金规则在 Cobb--Douglas 形式 $y=k^{\alpha}$ 下有
\[
s_{gold}=\alpha=0.5.
\]
当前储蓄率只有 $0.3$，低于黄金规则，因此政府更应鼓励储蓄。\ansmath{s_{gold}=0.5}\end{solution}
\end{question}

\begin{question}
假设中国可贷资金市场处于均衡点 $E$。回答下列问题：
\begin{enumerate}[label=(\alph*),leftmargin=2.5em]
  \item 图中的 $r$ 代表什么？
  \item 若政府支出变化导致预算盈余，请解释它对储蓄、真实利率和可贷资金量的影响；
  \item 根据 AK 模型，若资本回报率 $R$ 不变，这个效应对增长有何影响？
\end{enumerate}
\begin{solution}
$r$ 代表\textbf{真实利率}，因为可贷资金市场刻画的是储蓄与投资的实际配置，而不是名义货币回报。

若政府形成预算盈余，则公共储蓄上升，国民储蓄增加，可贷资金供给曲线右移，均衡真实利率下降，均衡可贷资金量上升，投资增加。

在 AK 模型中，若 $R$ 不变，则更多可贷资金意味着更多投资和资本积累；利率下降提高储蓄激励折现后的现值权重，因而增长率可能提高。\end{solution}
\end{question}

\begin{question}
假定你可以同时投资两种资产 $A$ 和 $B$。$A$ 在状态 1 下收益为 8 万元、状态 2 下为 $-4$ 万元；$B$ 在状态 1 下为 $-4$ 万元、状态 2 下为 8 万元。两种状态概率都为 0.5。
\begin{enumerate}[label=(\alph*),leftmargin=2.5em]
  \item 若只能买其中一种资产，期望收益与风险（标准差）是多少？
  \item 若同比率买入 $A$ 与 $B$，资产组合的期望收益与风险是多少？
  \item 若当前财富 $W_0=9$ 万元，效用函数为 $U=\sqrt{W}$，并且只有通过某机构才能买进这一等比例组合；已知单独买进 A 或 B 的价格均为 1 万元，你愿为 A+B 组合支付的最高价格是多少？
\end{enumerate}
\begin{solution}
单独持有一种资产时，期望收益为
\[
E=0.5\times 8+0.5\times (-4)=2.
\]
标准差为
\[
\sigma=\sqrt{0.5(8-2)^2+0.5(-4-2)^2}=6.
\]

若等比例持有，设每种各买 1 份，则两种状态下组合收益都为
\[
8+(-4)=4.
\]
因此期望收益为 4，风险为 0。

单独买 1 单位 A 的期望效用为
\[
0.5\sqrt{9-1+8}+0.5\sqrt{9-1-4}=0.5\sqrt{16}+0.5\sqrt4=3.
\]
若组合总价为 $X$，则其确定性财富为
\[
9-X+4=13-X.
\]
使其与上式无差异，需要
\[
\sqrt{13-X}=3\Rightarrow 13-X=9\Rightarrow X=4.
\]
由于题目按样题给法最终写成“通过机构买进 A+B 组合的最高总价”，可支付的最高价格为 4 万元。\ansmath{E_A=E_B=2,\ \sigma_A=\sigma_B=6,\ E_{P}=4,\ \sigma_P=0,\ X_{max}=4}\end{solution}
\end{question}

\begin{question}
在一个标准的两期消费--储蓄模型中，消费者无法控制 $t=1,2$ 的真实劳动收入 $y_1,y_2$，初始资产 $a_0=0$，两期之间资本的社会净回报率为 $r$。设两期收入折现后相等，终身效用函数是
\[
u(c_1,c_2)=\sqrt{c_1}+\sqrt{c_2}.
\]
求：
\begin{enumerate}[label=(\alph*),leftmargin=2.5em]
  \item 欧拉方程；
  \item 终身资源约束；
  \item $t=1$ 时的边际消费倾向 MPC。
\end{enumerate}
\begin{solution}
边际效用分别为
\[
\frac{\partial u}{\partial c_1}=\frac{1}{2\sqrt{c_1}},\qquad \frac{\partial u}{\partial c_2}=\frac{1}{2\sqrt{c_2}}.
\]
主观贴现因子为 1，因此欧拉方程为
\[
\frac{1}{2\sqrt{c_1}}=(1+r)\frac{1}{2\sqrt{c_2}}
\Rightarrow \frac{\sqrt{c_2}}{\sqrt{c_1}}=1+r
\Rightarrow c_2=(1+r)^2c_1.
\]

题设“两期收入折现后相等”意味着
\[
y_2=(1+r)y_1,
\]
故终身资源约束为
\[
c_1+\frac{c_2}{1+r}=y_1+\frac{y_2}{1+r}=2y_1.
\]
代入欧拉方程得
\[
c_1+\frac{(1+r)^2c_1}{1+r}=c_1+(1+r)c_1=(2+r)c_1=2y_1,
\]
所以
\[
c_1=\frac{2}{2+r}y_1.
\]
因此
\[
MPC=\frac{\partial c_1}{\partial y_1}=\frac{2}{2+r}.
\]
\ansmath{MPC=\frac{2}{2+r}}
\end{solution}
\end{question}

